\begin{table}
\centering
\caption{Dark matter fraction of rotationally-supported KURVS-CDFS galaxies.}
\label{tab3:fdm}
\begin{tabular}{cc} 
\hline
\hline
KURVS ID & $f_{\rm DM}(\leqslant R_{\rm eff})$  \\
(1) & (2) \\
\hline
11 & 0.42 $\pm$ 0.09 \\
13 & 0.62 $\pm$ 0.07 \\
15 & 0.31 $\pm$ 0.13 \\
16 & 0.50 $\pm$ 0.08 \\
17 & 0.69 $\pm$ 0.06 \\
21$^{\star}$ & 0.0  \\
3 & 0.19 $\pm$ 0.13 \\
7 & 0.46 $\pm$ 0.11 \\
8 & 0.67 $\pm$ 0.06 \\
9 & 0.54 $\pm$ 0.07 \\
\hline
\hline
\end{tabular}
\\[2mm] %You can adjust how far below the table the text should appear
\begin{flushleft}
{\bf Note.} 
(1) KURVS-ID. 
(2) Dark matter fraction within the effective radius and $1 \sigma$ uncertainties. This is obtained on seeing- and pressure support- corrected rotation curves using equation~\ref{fdm}. Galaxies are sorted in order of decreasing $v_{\rm rot}/\sigma_0$. 
$^{\star}$Galaxy KURVS-21 has a highly asymmetric velocity field, that makes it difficult to find a best-fitting model with \textsc{galpak$^{\rm 3D}$}. 
While we report its dark matter fraction for completeness, this measurements is highly unreliable and we flag this galaxy in the relevant plots.
\end{flushleft}
\end{table}
